% GENERATED FILE. Edit canonical lessons, then run npm run generate:books.
% canonical-source-hash: fnv1a64:d62694654c65f87c
% canonical-lessons: KA-C240-bogalu, KA-C240-kugu, KA-C240-mulugu, KA-C240-vadisu, KA-C240-kalpisiko

\chapter{To bark, To shout, To sink, To argue, To imagine}
\label{ch:ka-to-bark-to-shout-to-sink-to-argue-to-imagine}
\hlchaptermodality{240}

\begin{chapteropening}
\textbf{By the end of this chapter:} \emph{I can say to bark, to shout, to sink, to argue and to imagine.}

Complete the last lesson of chapter 240: I can say to bark, to shout, to sink, to argue and to imagine.
\end{chapteropening}

\section[\texorpdfstring{\kn{ಬೊಗಳು} (bogaḷu)}{ಬೊಗಳು (bogaḷu)}]{\kn{ಬೊಗಳು} (bogaḷu) — to bark}

\label{lesson:KA-C240-bogalu}



\begin{quote}
Before the new one: say the Kannada for to pay a visit, then the Kannada for to hug.
\end{quote}

\subsection*{You'll want to know: \kn{ಬೊಗಳು}}
\textbf{\kn{ಬೊಗಳು}} — \emph{bogaḷu} — \textquotedblleft{}to bark\textquotedblright{}.

\begin{grammarlens}[title={What you've built}]
One more word for this chapter.
\end{grammarlens}

\subsection*{Your turn}
\begin{itemize}
\raggedright
  \item \emph{Say it:} \emph{bogaḷu}
  \item \emph{Say it:} \emph{bogaḷu}, once more
  \item \emph{Say it:} say \emph{bogaḷu}
  \item \emph{Recall:} say the Kannada for to pay a visit, then the Kannada for to hug, then say \emph{bogaḷu} again
\end{itemize}

\begin{tcolorbox}[breakable,colback=teal!4,colframe=teal!35!black,title={Before you move on}]
What does \kn{ಬೊಗಳು} mean? (\textquotedblleft{}To bark\textquotedblright{}.) Say it once more.
\end{tcolorbox}
\section[\texorpdfstring{\kn{ಕೂಗು} (kūgu)}{ಕೂಗು (kūgu)}]{\kn{ಕೂಗು} (kūgu) — to shout, to cry out}

\label{lesson:KA-C240-kugu}



\begin{quote}
Before the new one: say the Kannada for to hug, then the Kannada for to bark.
\end{quote}

\subsection*{You'll want to know: \kn{ಕೂಗು}}
\textbf{\kn{ಕೂಗು}} — \emph{kūgu} — \textquotedblleft{}to shout, to cry out\textquotedblright{}.

\begin{grammarlens}[title={What you've built}]
One more word for this chapter.
\end{grammarlens}

\subsection*{Your turn}
\begin{itemize}
\raggedright
  \item \emph{Say it:} \emph{kūgu}
  \item \emph{Say it:} \emph{kūgu}, once more
  \item \emph{Say it:} say \emph{kūgu}
  \item \emph{Recall:} say the Kannada for to hug, then the Kannada for to bark, then say \emph{kūgu} again
\end{itemize}

\begin{tcolorbox}[breakable,colback=teal!4,colframe=teal!35!black,title={Before you move on}]
What does \kn{ಕೂಗು} mean? (\textquotedblleft{}To shout, to cry out\textquotedblright{}.) Say it once more.
\end{tcolorbox}
\section[\texorpdfstring{\kn{ಮುಳುಗು} (muḷugu)}{ಮುಳುಗು (muḷugu)}]{\kn{ಮುಳುಗು} (muḷugu) — to sink; to set (of the sun)}

\label{lesson:KA-C240-mulugu}



\begin{quote}
Before the new one: say the Kannada for to bark, then the Kannada for to shout.
\end{quote}

\subsection*{You'll want to know: \kn{ಮುಳುಗು}}
\textbf{\kn{ಮುಳುಗು}} — \emph{muḷugu} — \textquotedblleft{}to sink; to set (of the sun)\textquotedblright{}.

\begin{grammarlens}[title={What you've built}]
One more word for this chapter.
\end{grammarlens}

\subsection*{Your turn}
\begin{itemize}
\raggedright
  \item \emph{Say it:} \emph{muḷugu}
  \item \emph{Say it:} \emph{muḷugu}, once more
  \item \emph{Say it:} say \emph{muḷugu}
  \item \emph{Recall:} say the Kannada for to bark, then the Kannada for to shout, then say \emph{muḷugu} again
\end{itemize}

\begin{tcolorbox}[breakable,colback=teal!4,colframe=teal!35!black,title={Before you move on}]
What does \kn{ಮುಳುಗು} mean? (\textquotedblleft{}To sink; to set (of the sun)\textquotedblright{}.) Say it once more.
\end{tcolorbox}
\section[\texorpdfstring{\kn{ವಾದಿಸು} (vādisu)}{ವಾದಿಸು (vādisu)}]{\kn{ವಾದಿಸು} (vādisu) — to argue}

\label{lesson:KA-C240-vadisu}



\begin{quote}
Before the new one: say the Kannada for to shout, then the Kannada for to sink.
\end{quote}

\subsection*{You'll want to know: \kn{ವಾದಿಸು}}
\textbf{\kn{ವಾದಿಸು}} — \emph{vādisu} — \textquotedblleft{}to argue\textquotedblright{}.

\begin{grammarlens}[title={What you've built}]
One more word for this chapter.
\end{grammarlens}

\subsection*{Your turn}
\begin{itemize}
\raggedright
  \item \emph{Say it:} \emph{vādisu}
  \item \emph{Say it:} \emph{vādisu}, once more
  \item \emph{Say it:} say \emph{vādisu}
  \item \emph{Recall:} say the Kannada for to shout, then the Kannada for to sink, then say \emph{vādisu} again
\end{itemize}

\begin{tcolorbox}[breakable,colback=teal!4,colframe=teal!35!black,title={Before you move on}]
What does \kn{ವಾದಿಸು} mean? (\textquotedblleft{}To argue\textquotedblright{}.) Say it once more.
\end{tcolorbox}
\section[\texorpdfstring{\kn{ಕಲ್ಪಿಸಿಕೊ} (kalpisiko)}{ಕಲ್ಪಿಸಿಕೊ (kalpisiko)}]{\kn{ಕಲ್ಪಿಸಿಕೊ} (kalpisiko) — to imagine}

\label{lesson:KA-C240-kalpisiko}



\begin{quote}
Before the new one: say the Kannada for to sink, then the Kannada for to argue.
\end{quote}

\subsection*{You'll want to know: \kn{ಕಲ್ಪಿಸಿಕೊ}}
\textbf{\kn{ಕಲ್ಪಿಸಿಕೊ}} — \emph{kalpisiko} — \textquotedblleft{}to imagine\textquotedblright{}.

\begin{grammarlens}[title={What you've built}]
One more word for this chapter.
\end{grammarlens}

\subsection*{Your turn}
\begin{itemize}
\raggedright
  \item \emph{Say it:} \emph{kalpisiko}
  \item \emph{Say it:} \emph{kalpisiko}, once more
  \item \emph{Say it:} say \emph{kalpisiko}
  \item \emph{Recall:} say the Kannada for to sink, then the Kannada for to argue, then say \emph{kalpisiko} again
\end{itemize}

\begin{tcolorbox}[breakable,colback=teal!4,colframe=teal!35!black,title={Before you move on}]
What does \kn{ಕಲ್ಪಿಸಿಕೊ} mean? (\textquotedblleft{}To imagine\textquotedblright{}.) Say it once more.
\end{tcolorbox}
